\section{Identifiability, Active Repair, and Rational Unresolvedness}

Gates 110--115 expose the informational condition hidden inside Gate 109's repair oracle. A controller may terminate perfectly as a state machine and still be unable to determine which classical verdict is true. The issue is not computational weakness but observational aliasing: different hidden worlds may produce the same evidence while requiring different answers.

\subsection{Aliasing as a limit on truth-guaranteeing repair}

Gate 110 gives two hidden worlds with the same visible observation but different target verdicts. Any policy that depends only on that observation must return the same output in both worlds. Uniform correctness would require different outputs. Therefore no observation-only policy can be truth-guaranteeing for both worlds.

The result yields a precise separation:
\[
\text{repair termination}
\not\Rightarrow
\text{truth}.
\]
A Gate-109-style controller can reach a classical state and nevertheless be wrong because the evidence channel never contained the distinction needed to choose the correct classical value. This impossibility is channel-relative. It does not say that the worlds are indistinguishable under every possible intervention.

\subsection{Active identification restores a truth guarantee}

Gate 111 therefore gives the agent a choice of experiment. A passive experiment preserves the Gate-110 alias. A supplied \texttt{disambiguate} experiment returns different signals in the two hidden worlds. The decoder maps those signals to the corresponding $T/F$ target, and the repaired nested state is both classically stable and truth-aligned.

The result changes the role of epistemic agency. The agent need not passively accept the information partition generated by its current sensor. It can act to refine that partition:
\[
\text{alias}
\longrightarrow
\text{intervention}
\longrightarrow
\text{discriminating observation}
\longrightarrow
\text{truth-aligned repair}.
\]
The separating experiment is supplied as part of the witness. Gate 111 therefore establishes a possibility result, not the claim that every alias admits such an intervention.

\subsection{The price of identifiability}

Gate 112 asks whether a perfect separating experiment should always be purchased. Under the symmetric two-world witness, a blind classical commitment has expected value $1/2$. Perfect diagnosis has gross value $1$ and net value $1-c$ at diagnostic cost $c$. Hence diagnosis strictly dominates blind commitment exactly when
\[
c<\frac12.
\]
Identifiability may therefore be available and nevertheless not be worth its price.

This cost threshold is specific to the encoded utilities. The structural result is broader: epistemic access and rational acquisition are separate questions.

\subsection{Rationally preserving the gap}

Gate 113 adds a third action that neither guesses nor purchases perfect information. The agent may remain unresolved. Safe deferral is assigned value
\[
\frac35,
\]
which exceeds blind commitment's value $1/2$. Perfect identification is still preferable when sufficiently cheap, but it beats unresolvedness exactly when
\[
c<\frac25.
\]
At the high-cost witness $c=1/2$, the controller chooses \texttt{remainUnresolved}, preserving a world status $N$ under a trusted model status $T$.

Thus there exist coherent finite value models in which the rational response to an epistemic gap is neither forced classicalization nor further inquiry:
\[
\boxed{N\text{ may be worth preserving.}}
\]
The numerical value $3/5$ is an explicit modelling choice. The theorem establishes the existence of such a rational-control regime, not a universal utility of ignorance.

\subsection{A general identifiability criterion}

Gate 114 abstracts away from the finite two-world examples. Let
\[
\mathsf{observe}:W\to O,
\qquad
\mathsf{target}:W\to V.
\]
A decoder $d:O\to V$ is truth-guaranteeing when
\[
\forall w\in W,\qquad d(\mathsf{observe}(w))=\mathsf{target}(w).
\]
The verified theorem states
\[
\exists d\;\forall w,\ d(\mathsf{observe}(w))=\mathsf{target}(w)
\quad\Longleftrightarrow\quad
\forall w_1,w_2,
\]
\[
\mathsf{observe}(w_1)=\mathsf{observe}(w_2)
\Rightarrow
\mathsf{target}(w_1)=\mathsf{target}(w_2).
\]
In words, a truth-guaranteeing decoder exists exactly when the target is constant on every observation fiber.

Gate 110 is a direct violation of this condition. Gate 111's disambiguating experiment satisfies it. The reverse direction constructs a generic decoder using classical choice and an arbitrary default outside the image of the observation map. The theorem is therefore existential, not an efficient decoder-synthesis algorithm.

A central consequence is task sensitivity. Full hidden-world identification is stronger than necessary. Two worlds may remain observationally indistinguishable provided the requested target gives them the same answer.

\subsection{Minimal target-relevant experimentation}

Gate 115 turns this consequence into an experiment-design witness. There are four hidden worlds $a,b,c,d$ with target values
\[
T,\ T,\ F,\ B.
\]
Two binary experiments are available. Neither single experiment, nor the empty family, satisfies Gate 114's target-aware fiber condition. The pair does. Yet even the pair fails to identify the hidden world completely because $a$ and $b$ remain observationally aliased.

That remaining alias is harmless for the specified task because
\[
\mathsf{target}(a)=\mathsf{target}(b)=T.
\]
Thus the two-experiment family is inclusion-minimal inside the explicit menu for target identification while remaining insufficient for full-state reconstruction.

This yields a more economical conception of epistemic adequacy:
\begin{quote}
The agent need not learn every distinction in the world. It needs to learn the distinctions across which the correct answer changes.
\end{quote}
The next step is to remove one last idealization. Gates 110--115 use deterministic observation partitions. Gate 116 allows the same signal to occur in target-different worlds with different probabilities, replacing perfect identification by fallible evidence and posterior confidence.
